\begin{beginnerstatement}[Kurz erklärt: FastAPI und HTTP]

Das Backend ist der Teil im Hintergrund; hier wird es mit der Sprache Python
und dem Werkzeug FastAPI erstellt. HTTP legt Regeln dafür fest, wie Programme
Anfragen und Antworten über ein Netzwerk übertragen. Eine API ist eine
vereinbarte Schnittstelle, über die das Frontend auf diesem Weg mit dem Backend
kommuniziert.
Eine Route oder ein Endpoint ist eine ansprechbare Adresse dieser API, wobei
\texttt{/api} die gemeinsamen Prüfpfade gruppiert. \texttt{/api/health} meldet
als Liveness-Prüfung, ob der Prozess lebt; \texttt{/api/ready} meldet als
Readiness-Prüfung, ob er einsatzbereit ist, und prüft bei aktivierter Datenbank
auch deren Verbindung. Ein HTTP-Statuscode fasst das Ergebnis als Zahl
zusammen; HTTP~503 bedeutet hier, dass der Dienst vorübergehend nicht bereit
ist. Ein Container verpackt einen Dienst mit seiner Laufzeit in eine getrennte
Einheit, und das FastAPI-Backend läuft beim Deployment in einem solchen
Container. Produktspezifische Dienste und eine Authentifizierung enthält diese
Baseline noch nicht.

\end{beginnerstatement}
